% Generated by thesis/ch7/exp2/completeness_table.py -- do not edit by hand.
\begin{table}
  \thesistablestyle
  \centering
  \footnotesize
  \begin{tabular}{@{}lccc@{}}
    \toprule
    & Adult & Churn & Heart-100 \\
    \midrule
    Logit $R^2$ on the signals, train: median (min) & 1.000 (0.998) & 0.992 (0.943) & 0.964 (0.936) \\
    Logit $R^2$, test: median (min) & 1.000 (0.999) & 0.989 (0.930) & 0.924 (0.849) \\
    Linear terms alone, $R^2$ train (median) & 0.993 & 0.671 & 0.915 \\
    Residual SD (logit), train: median (max) & 0.032 (0.242) & 0.158 (0.467) & 0.754 (1.070) \\
    Networks with residual SD $\geq 0.10$ & 1 of 10 & 9 of 10 & 10 of 10 \\
    Parts used (all) & 61 (94) & 66 (84) & 52 (60) \\
    Parts mapping to no signal & 0 & 0 & 1 \\
    Mapped parts with residual SD $\geq 0.10$ & 1 & 16 & 37 \\
    Candidate signals in the residual & 4 & 5 & 4 \\
    Complete by the criterion & no & no & no \\
    \bottomrule
  \end{tabular}
  \caption{Is each signal list complete? Each network's logit is projected on the data functions of its table's signals (training rows). Complete when every network's residual has an SD below 0.10 logit and the residual scan (one variable and variable pairs, permutation tested) finds no candidate signal. Source: \texttt{experiments/signal\_completeness}.}
  \label{tab:exp2-signal-completeness}
\end{table}
